\chaptertitle{Contrafactus}

                                  Contrafactus
   The Crab has invited a small group of friends over to watch the Saturday
   afternoon football game on television. Achilles has already arrived, but the
   Tortoise and his friend the Sloth are still awaited.


Achilles: Could that be our friends, a-riding up on that unusual one-wheeled vehicle?

  (The Sloth and Tortoise dismount and come in.)

Crab: Ah, my friends, I'm so glad you could make it. May I present my old and beloved
   acquaintance, Mr. Sloth-and this is Achilles. I believe you know the Tortoise.
Sloth: This is the first time I can recall making the acquaintance of a Bicyclops. Pleased
   to meet you, Achilles. I've heard many fine things said about the bicyclopean species.
Achilles: Likewise, I'm sure. May I ask about your elegant vehicle? Tortoise: Our tandem
   unicycle, you mean? Hardly elegant. It's just a way for two to get from A to B, at the
   same speed.
Sloth: It's built by a company that also makes teeter-teeters.
Achilles: I see, I see. What is that knob on it?
Sloth: That's the gearshift.
Achilles: Aha! And how many speeds does it have?
Tortoise: One, including reverse. Most models have fewer, but this is a special model.
Achilles: It looks like a very nice tandem unicycle. Oh, Mr. Crab, I wanted to tell you
   how much I enjoyed hearing your orchestra perform last night.
Crab: Thank you, Achilles. Were you there by any chance, Mr. Sloth? Sloth: No, I
   couldn't make it, I'm sad to say. I was participating in a mixed singles ping-ping
   tournament. It was quite exciting because my team was involved in a one-way tie for
   first place.
Achilles: Did you win anything?
Sloth: Certainly did-a two-sided Mobius strip made out of copper; it is silver-plated on
   one side, and gold-plated on the other. Crab: Congratulations, Mr. Sloth.
Sloth: Thank you. Well, do tell me about the concert.
Crab: It was a most enjoyable performance. We played some pieces by the Bach twins
Sloth: The famous Job and Sebastian?
Crab: One and the same. And there was one work that made me think of you, Mr. Sloth-a
   marvelous piano concerto for two left hands. The




Contrafactus                                                                            630

   next-to-last (and only) movement was a one-voice fugue. You can't imagine its
   intricacies. For our finale, we played Beethoven's Ninth Zenfunny. At the end,
   everyone in the audience rose and clapped with one hand. It was overwhelming.
Sloth: Oh, I'm sorry I missed it. But do you suppose it's been recorded: At home I have a
   fine hi-fi to play it on-the best two-channel monaural system money can buy.
Crab: I'm sure you can find it somewhere. Well, my friends, the game is about to begin.
Achilles: Who is playing today, Mr. Crab?
Crab: I believe it's Home Team versus Visitors. Oh, no-that was last week. I think this
   week it's Out-of-Towners.
Achilles: I'm rooting for Home Team. I always do.
Sloth: Oh, how conventional. I never root for Home Team. The closer a team lives to the
   antipodes, the more I root for it.
Achilles: Oh, so you live in the Antipodes? I've heard it's charming to live there, but I
   wouldn't want to visit them. They're so far away.
Sloth: And the strange thing about them is that they don't get any closer no matter which
   way you travel.
Tortoise: That's my kind of place.
Crab: It's game time. I think I'll turn on the TV.

   (He walks over to an enormous cabinet with a screen, underneath which is an
   instrument panel as complicated as that of a jet airplane. He flicks a knob, and the
   football stadium a ears in bright vivid color on the screen.)

Announcer: Good afternoon, fans. Well, it looks like that time of year has rolled around
   again when Home Team and Out-of-Town face each other on the gridiron and play out
   their classic pigskin rivalry. It's been drizzling on and off this afternoon, and the field's
   a little wet, but despite the weather it promises to be a fine game, especially with that
   GREAT pair of eighth-backs playing for Home Team, Tedzilliger and Palindromi.
   And now, here's Pilipik, kicking off for Home Team. It's in the air! Flampson takes it
   for Out-of-Towners, and runs it back he's to the 20, the 25, the 30, and down at the 32.
   That was Mool in on the tackle for Home Team.
Crab: A superb runback! Did you see how he was ALMOST tackled by Quilker-but
   somehow broke away?
Sloth: Oh, don't be silly, Crab. Nothing of the kind happened. Quilker did NOT tackle
   Flampson. There's no need to confuse poor Achilles (or the rest of us) with hocus-
   pocus about what "almost" happened. It's a fact-with no "almost" 's, "if "'s, "and" 's, or
   "but" 's.
Announcer: Here's the instant replay. Just watch number 79, Quilker, come in from the
   side, surprising Flampson, and just about tackle him!
Sloth: "Just about"! Bah!
'Achilles: Such a graceful maneuver! What would we do without instant replays?




Contrafactus                                                                                631

Announcer: It's first down and 10 for Out-of-Town. Noddle takes the ball, hands off to
   Orwix-it's a reverse-Orwix runs around to the right, handing off to Flampson-a double
   reverse, folks!-and now
Flampson hands it to Treefig, who's downed twelve yards behind scrimmage. A twelve-
   yard loss on a triple reverse!
Sloth: I love it! A sensational play!
Achilles: But, Mr. S, I thought you were rooting for Out-of-Town. They lost twelve yards
   on the play.
Moth: They did? Oh, well-who cares, as long as it was a beautiful play?
Let's see it again.

   ( ... and so the first half of the game passes. Towards the end of the third quarter, a
   particularly crucial play comes up for Home Team. They are behind by eight points.
   It's third down and 10, and they badly need a first down.)

Announcer: The ball is hiked to Tedzilliger, who fades back, looking-for a receiver, and
  fakes to Quilker. There's Palindromi, playing wide right, with nobody near him.
  Tedzilliger spots him and fires a low pass to him. Palindromi snatches it out of the air,
  and- (There is an audible groan from the crowd.)--oh, he steps out of bounds! What a
  crushing blow for Home Team, folks! If Palindromi hadn't stepped out of bounds, he
  could've run all the way to the end zone for a touchdown!
Let's watch the subjunctive instant replay.

   (And on the screen the same lineup appears as before.)

   The ball is hiked to Tedzilliger, who fades back, looking for a receiver, and fakes to
   Quilker. There's Palindromi, playing wide right, with nobody near him. Tedzilliger
   spots him, and fires a low pass to him. Palindromi snatches it out of the air, and-
   (There is an audible gasp from the crowd.)-he almost steps out of bounds! But he's
   still in bounds, and it's clear all the way to the end zone! Palindromi streaks in, for a
   touchdown for Home Team! (The stadium breaks into a giant roar of approval.) Well,
   folks, that's what would've happened if Palindromi hadn't stepped out of bounds.

Achilles: Wait a minute ... WAS there a touchdown, or WASN'T there?
Crab: Oh, no. That was just the subjunctive instant replay. They simply followed a
   hypothetical a little way out, you know.
Sloth: That is the most ridiculous thing I ever heard of! Next thing you know, they'll be
   inventing concrete earmuffs.
Tortoise: Subjunctive instant replays are a little unusual, aren't they?
Crab: Not particularly, if you have a Subjunc-TV.
Achilles: Is that one grade below a junk TV?
Crab: Not at all! It's a new kind of TV, which can go into the subjunctive mode. They're
   particularly good for football games and such. I just got mine.
Achilles: Why does it have so many knobs and fancy dials?




Contrafactus                                                                             632

Crab: So that you can tune it to the proper channel. There are many channels
   broadcasting in the subjunctive mode, and you want to be able to select from them
   easily.
Achilles: Could you show us what you mean? I'm afraid I don't quite understand what all
   this talk of "broadcasting in the subjunctive mode" is about.
Crab: Oh, it's quite simple, really. You can figure it out yourself. I'm going into the
   kitchen to fix some French fries, which I know are Mr. Sloth's weakness.
Sloth: Mmmmm! Go to it, Crab! French fries are my favorite food. Crab: What about the
   rest of you?
Tortoise: I could devour a few.
Achilles: Likewise. But wait-before you go into the kitchen, is there some trick to using
   your Subjunc-TV?
Crab: Not particularly. Just continue watching the game. and whenever there's a near
   miss of some sort, or whenever you wish things had gone differently in some way, just
   fiddle with the dials, and see what happens. You can't do it any harm, though you may
   pick up some exotic channels. (And he disappears into the kitchen.)
Achilles: I wonder what he means by that. Oh well, let's get back to this game. I was
   quite wrapped up in it.
Announcer: It's fourth down for Out-of-Town, with Home Team receiving. Out-of-Town
   is in punt formation, with Tedzilliger playing deep. Orwix is back to kick-and he gets
   a long high one away. It's coming down near Tedzilliger
Achilles: Grab it, Tedzilliger! Give those Out-of-Towners a run for their money!
Announcer: -and lands in a puddle-KERSPLOSH! It takes a weird bounce! Now Sprunk
   is madly scrambling for the ball! It looks like it just barely grazed Tedzilliger on the
   bounce, and then slipped away from himit's ruled a fumble. The referee is signaling
   that the formidable Sprunk has recovered for Out-of-Town on the Home Team 7! It's a
   bad break for Home Team. Oh, well, that's the way the cookie crumbles.
Achilles: Oh, no! If only it hadn't been raining ... (Wrings his hands in despair.)
Sloth: ANOTHER of those confounded hypotheticals! Why are the rest of you always
   running off into your absurd worlds of fantasy? If I were you, I would stay firmly
   grounded in reality. "No subjunctive nonsense" is my motto. And I wouldn't abandon
   it even if someone offered me a hundred-nay, a hundred and twelve-French fries.
Achilles: Say, that gives me an idea. Maybe by suitably fiddling with these knobs, I can
   conjure up a subjunctive instant replay in which it isn't raining, there's no puddle, no
   weird bounce, and Tedzilliger doesn't
fumble. I wonder ... (Walks up to the Subjunc-TV and stares at it.) But I haven't any idea
   what these different knobs do. (Spins a few at random.)
Announcer: It's fourth down for Out-of-Town, with Home Team receiv-




Contrafactus                                                                           633

  ing. Out-of-Town is in punt formation, with Tedzilliger playing deep. Orwix is back to
  kick-and he gets a long high one away. It's coming down near Tedzilliger
Achilles: Grab it, Tedzilliger! Give those Out-of-Towners a run for their money!
Announcer: -and lands in a puddle-KERSPLOSH! Oh-it bounces right into his arms!
  Now Sprunk is madly scrambling after him, but he's got - good blocking, and he steers
  his way clear of the formidable Sprunk, and now he's got an open field ahead of him.
  Look at that, folks! He's to the 50, the 40, the 30, the 20, the 10-touchdown, Home
  Team! (Huge cheers from the Home Team side.) Well, fans, that's how it would have
  gone, if footballs were spheres instead of oblate spheroids! But in reality, Home Team
  loses the ball, and Out-of-Towners take over on the Home Team 7-yard line. Oh, well,
  that's the way the ball bounces.
Achilles: What do you think of THAT, Mr. Sloth?

   (And Achilles gives a smirk in the direction of the Sloth, but the latter is completely
   oblivious to its devastating effect, as he is busy watching ,the Crab arrive with, a large
   platter with a hundred and twelve-nay, a hundred-large and delicious French fries,
   and napkins for all.)

Crab: So how do you three find my Subjunc-TV?
Sloth: Most disappointing, Crab, to be quite frank. It seems to be badly out of order. It
   makes pointless excursions into nonsense at least half the time. If it belonged to me, I
   would give it away immediately to someone like you, Crab. But of course it doesn't
   belong to me.
Achilles: It's quite a strange device. I tried to rerun a play to see how it would have gone
   under different weather conditions, but the thing seems to have a will of its own!
   Instead of changing the weather, it changed the football shape to ROUND instead of
   FOOTBALL-SHAPED! Now tell me-how can a football not be shaped like a football?
   That's a contradiction in terms. How preposterous!
Crab: Such tame games! I thought you'd surely find more interesting subjunctives. How
   would you like to see how the last play would have looked if the game had been
   baseball instead of football?
   Tortoise: Oh! An outstanding idea!

(The Crab twiddles two knobs, and steps back.)

Announcer: There are four away, and---
Achilles: FOUR away!?
Announcer: That's right, fans-four away. When you turn football into baseball,
  SOMETHING'S got to give! Now as I was about to say, there are four away, with
  Out-of-Town in the field, and Home Team up. Tedzilliger is at bat. Out-of-Town is in
  bunt formation. Orwix raises his arm to pitch-and he gets a long high ball away. It's
  heading straight for Tedzilliger
Achilles: Smash it, Tedzilliger! Give those Out-of-Towners a home run for their money!




Contrafactus                                                                             634

Announcer: -but it seems to be a spitball, as it takes a strange curve. Now Sprunk is
   madly scrambling for the ball! It looks like it just barely grazed Tedzilliger's bat, then
   bounced off it-it's ruled a fly ball. The umpire is signaling that the formidable Sprunk
   has caught it for Out-of-Town, to end the seventh inning. It's a bad break for Home
   Team. That's how the last play would have looked, football fans, if this had been a
   game of baseball.
Sloth: Bah! You might as well transport this game to the Moon.
Crab: No sooner said than done! Just a twiddle here, a twiddle there ...

   (On the screen there appears a desolate crater-pitted field, with two teams in space
   suits facing each other, immobile. All at once, the two teams fly into motion, and the
   players are making great bounds into the air, sometimes over the heads of other
   players. The ball is thrown into the air, and sails so high that it almost disappears,
   and then slowly comes floating down into the arms of one space-suited player, roughly
   a quarter-mile from where it was released.)

Announcer: And there, friends, you have the subjunctive instant replay as it would have
   happened on the Moon. We'll be right back after this important commercial message
   from the friendly folks who brew Glumpf Beer-my favorite kind of beer!
Sloth: If I weren't so lazy, I would take that broken TV back to the dealer myself! But
   alas, it's my fate to be a lazy Sloth ... (Helps himself to a large gob of French fries.)
Tortoise: That's a marvelous invention, Mr. Crab. May I suggest a hypothetical?
Crab: Of course!
Tortoise: What would that last play have looked like if space were four-dimensional?
Crab: Oh, that's a complicated one, Mr. T, but I believe I can code it into the dials. Just a
   moment.

   (He steps up, and, for the first time, appears to be using the full power of the control
   panel of his Subjunc-TV, turning almost every knob two or three times, and carefully
   checking various meters. Then he steps back with a satisfied expression on his face.)

I think this should do it.

Announcer: And now let's watch the subjunctive instant replay.

   (A confusing array of twisted pipes appears on the screen. It grows larger, then
   smaller, and for a moment seems to do something akin to rotation. Then it turns into a
   strange mushroom-shaped object, and back to a bunch of pipes. As it metamorphoses
   from this into other bizarre shapes, the announcer gives his commentary.)

Tedzilliger's fading back to pass. He spots Palindromi ten yards outfield, and passes it to
  the right and outwards-it looks good! Palindromi's at the 35-yard plane, the 40, and
  he's tackled on his own




Contrafactus                                                                              635

   43-yard plane. And there you nave it, 3-L tans, as it would have looked if football
   were played in four spatial dimensions.
Achilles: What is it you are doing, Mr. Crab, when you twirl these various dials on the
   control panel?
Crab: I am selecting the proper subjunctive channel. You see, there are all sorts of
   subjunctive channels broadcasting simultaneously, and I want to tune in precisely that
   one which represents the kind of hypothetical which has been suggested.
Achilles: Can you do this on any TV?
Crab: No, most TV's can't receive subjunctive channels. They require a special kind of
   circuit which is quite difficult to make.
Sloth: How do you know which channel is broadcasting what? Do you look it up in the
   newspaper?
Crab: I don't need to know the channel's call letters. Instead, I tune it in by coding, in
   these dials, the hypothetical situation which I want to be represented. Technically, this
   is called "addressing a channel by its counterfactual parameters". There are always a
   large number of channels broadcasting every conceivable world. All the channels
   which carry worlds that are "near" to each other have call letters that are near
to each other, too.
Tortoise: Why did you not have to turn the dials at all, the first time we saw a subjunctive
   instant replay?
Crab: Oh, that was because I was tuned in to a channel which is very near to the Reality
   Channel, but ever so slightly off. So every once in a while, it deviates from reality. It's
   nearly impossible to tune EXACTLY into the Reality Channel. But that's all right,
   because it's so dull. All their instant replays are straight! Can you imagine? What a
   bore!
Sloth: I find the whole idea of Subjunc-TV's one giant bore. But perhaps I could change
   my mind, if I had some evidence that your machine here could handle an
   INTERESTING counterfactual. For example, how would that last play have looked if
   addition were not commutative?
Crab: Oh me, oh my! That change is a little too radical, I'm afraid, for this model. I
   unfortunately don't have a Superjunc-TV, which is the top of the line. Superjunc-TV's
   can handle ANYTHING you throw at them.
Sloth: Bah!
Crab: But look-I can do ALMOST as well. Wouldn't you like to see how the last play
   would have happened if 13 were not a prime number? Sloth: No thanks! THAT
   doesn't make any sense! Anyway, if I were the last play, I'd be getting pretty tired of
   being trotted out time and again in new garb for the likes of you fuzzy-headed
   concept-slippers. Let's get on with the game!
Achilles: Where did you get this Subjunc-TV, Mr. Crab?
Crab: Believe it or not, Mr. Sloth and I went to a country fair the other evening, and it
   was offered as the first prize in a lottery. Normally I don't indulge in such frivolity, but
   some crazy impulse grabbed me, and I bought one ticket.




Contrafactus                                                                               636

Achilles: What about you Mr. Sloth?
Sloth: I admit, I bought one, just to humor old Crab.
Crab: And when the winning number was announced, I found, to my amazement, that I'd
   won the lottery!
Achilles: Fantastic! I've never known anyone who won anything in a lottery before!
Crab: I was flabbergasted at my good fortune.
Sloth: Don't you have something else to tell us about that lottery, Crab?
Crab: Oh, nothing much. It's just that my ticket number was 129. Now when they
   announced the winning number, it was 128 just one off. Sloth: So you see, he actually
   didn't win it at all. Achilles: He ALMOST won, though ...
Crab: I prefer to say that I won it, you see. For I came so terribly close . . If my number
   had been only one smaller, I would have won. Sloth: But unfortunately, Crab, a miss is
   as good as a mile.
Tortoise: Or as bad. What about you, Mr. Sloth? What was your number:
Sloth: Mine was 256-the next power of 2 above 128. Surely, that counts as a hit, if
   anything does! I can't understand why, however, those fair officials-those UNfair
   officials-were so thickheaded about it. They refused to award me my fully deserved
   prize. Some other joker claimed HE deserved it, because his number was 128. 1 think
   my number was far closer than His, but you can't fight City Hall.
Achilles: I'm all confused. If you didn't win the Subjunc-TV after all, Mr. Crab, then how
   can we have been sitting here all afternoon watching it? It seems as if we ourselves
   have been living in some sort of hypothetical world that would have been, had
   circumstances just been ever so slightly different ...
Announcer: And that, folks, was how the afternoon at Mr. Crab's would have been spent,
   had he won the Subjunc-TV. But since he didn't, the four friends simply spent a
   pleasant afternoon watching Home Team get creamed, 128-0. Or was it 256-0? Oh
   well, it hardly matters, in five-dimensional Plutonian steam hockey.




Contrafactus                                                                           637


                          Artificial Intelligence:
                                 Prospects
                     "Almost" Situations and Subjunctives
AFTER READING Contrafactus, a friend said to me, "My uncle was almost President of
the U.S.!" "Really?" I said. "Sure," he replied, "he was skipper of the PT 108." (John F.
Kennedy was skipper of the PT 109.)
         That is what Contrafactus is all about. In everyday thought, we are constantly
manufacturing mental variants on situations we face, ideas we have, or events that
happen, and we let some features stay exactly the same while others "slip". What features
do we let slip? What ones do we not even consider letting slip? What events are
perceived on some deep intuitive level as being close relatives of ones which really
happened? What do we think "almost" happened or "could have" happened, even though
it unambiguously did not? What alternative versions of events pop without any conscious
thought into our minds when we hear a story? Why do some counterfactuals strike us as
"less counterfactual" than other counterfactuals? After all, it is obvious that anything that
didn't happen didn't happen. There aren't degrees of "didn't-happen-ness". And the same
goes for "almost" situations. There are times when one plaintively says, "It almost
happened", and other times when one says the same thing, full of relief. But the "almost"
lies in the mind, not in the external facts.
         Driving down a country road, you run into a swarm of bees. You don't just duly
take note of it; the whole situation is immediately placed in perspective by a swarm of
"replays" that crowd into your mind. Typically, you think, "Sure am lucky my window
wasn't open!"-or worse, the reverse: "Too bad my window wasn't closed!" "Lucky I
wasn't on my bike!" "Too bad I didn't come along five seconds earlier." Strange but-
possible replays: "If that had been a deer, I could have been killed!" "I bet those bees
would have rather had a collision with a rosebush." Even stranger replays: "Too bad
those bees weren't dollar bills!" "Lucky those bees weren't made of cement!" "Too bad it
wasn't just one bee instead of a swarm." "Lucky I wasn't the swarm instead of being me."
What slips naturally and what doesn't-and why?
         In a recent issue of The New Yorker magazine, the following excerpt from the
"Philadelphia Welcomat" was reprinted:'

   If Leonardo da Vinci had been born a female the ceiling of the Sistine Chapel
   might never have been painted.1




Artificial Intelligence:Prospects                                                        638

The New The New Yorker commented:

   And if Michelangelo had been Siamese twins, the work would have been
   completed in half the time.

The point of The New Yorker's comment is not that such counterfactuals are false; it is
more that anyone who would entertain such an idea-anyone who would "slip" the sex or
number of a given human being-would have to be a little loony. Ironically, though, in the
same issue, the following sentence, concluding a book review, was printed without
blushing:

   I think he [Professor Philipp Frank would have enjoyed both of these books
   enormously.2

Now poor Professor Frank is dead; and clearly it is nonsense to suggest that someone
could read books written after his death. So why wasn't this serious sentence also scoffed
at? Somehow, in some difficult-to-pin-down sense, the parameters slipped in this
sentence do not violate our sense of "possibility" as much as in the earlier examples.
Something allows us to imagine "all other things being equal" better in this one than in
the others. But why? What is it about the way we classify events and people that makes
us know deep down what is "sensible" to slip, and what is "silly":
        Consider how natural it feels to slip from the valueless declarative "I don't know
Russian" to the more charged conditional "I would like to know Russian" to the
emotional subjunctive "I wish I knew Russian" and finally to the rich counterfactual "If I
knew Russian, I would read Chekhov and Lermontov in the original". How flat and dead
would be a mind that saw nothing in a negation but an opaque barrier! A live mind can
see a window onto a world of possibilities.
        I believe that "almost" situations and unconsciously manufactured subjunctives
represent some of the richest potential sources of insight into how human beings organize
and categorize their perceptions of the world.
An eloquent co-proponent of this view is the linguist and translator George
Steiner, who, in his book After Babel, has written:

   Hypotheticals, 'imaginaries', conditionals, the syntax of counter-factuality and
   contingency may well be the generative centres of human speech.... [They] do more
   than occasion philosophical and grammatical perplexity. No less than future tenses
   to which they are, one feels, related, and with which they ought probably to be
   classed in the larger set of 'suppositionals' or `alternates', these `if' propositions are
   fundamental to the dynamics of human feeling... .

   Ours is the ability, the need, to gainsay or 'un-say' the world, to image and speak it
   otherwise.... We need a word which will designate the power, the compulsion of
   language to posit 'otherness'. . . . Perhaps 'alternity' will do: to define the `other than
   the case', the counter-factual propositions, images, shapes of will and evasion with
   which we charge our mental being and by means of which we build the changing,
   largely fictive milieu of our somatic and our social existence... .

Finally, Steiner sings a counterfactual hymn to counterfactuality:


Artificial Intelligence:Prospects                                                            639

It is unlikely that man, as we know him, would have survived without the fictive,
counter-factual, anti-determinist means of language, without the semantic capacity,
generated and stored in the `superfluous, zones of the cortex, to conceive of, to articulate
possibilities beyond the treadmill of organic decay and death .3

        The manufacture of "subjunctive worlds" happens so casually, -so naturally, that
we hardly notice what we are doing. We select from our fantasy a world which is close,
in some internal mental sense, to the real world. We compare what is real with what we
perceive as almost real. In so doing, what we gain is some intangible kind of perspective
on reality. The Sloth is a droll example of a variation on reality-a thinking being without
the ability to slip into subjunctives (or at least, who claims to be without the ability-but
you may have noticed that what he says is full of counterfactuals'.). Think how
immeasurably poorer our mental lives would be if we didn't have this creative capacity
for slipping out of the midst of reality into soft "what if'-s! And from the point of view of
studying human thought processes, this slippage is very interesting, for most of the time it
happens completely without conscious direction, which means that observation of what
kinds of things slip, versus what kinds don't, affords a good window on the unconscious
mind.
        One way to gain some perspective on the nature of this mental metric is to "fight
fire with fire". This is done in the Dialogue, where our "subjunctive ability" is asked to
imagine a world in which the very notion of
subjunctive ability is slipped, compared to what we expect. In the Dialogue, the first
subjunctive instant replay-that where Palindromi stays in bounds-is quite a normal thing
to imagine. In fact, it was inspired by a completely ordinary, casual remark made to me
by a person sitting next to me at a football game. For some reason it struck me and I
wondered what made it seem so natural to slip that particular thing, but not, say, the
number of the down, or the present score. From those thoughts, I went on to consider
other, probably less slippable features, such as the weather (that's in the Dialogue), the
kind of game (also in the Dialogue), and then even loonier variations (also in the
Dialogue). I noticed, though, that what was completely ludicrous to slip in one situation
could be quite slippable in another. For instance, sometimes you might spontaneously
wonder how things would be if the ball had a different shape (e.g., if you are playing
basketball with a half-inflated ball); other times that would never enter your mind (e.g.,
when watching a football game on TV).

Layers of Stability
It seemed to me then, and still does now, that the slippability of a feature of some event
(or circumstance) depends on a set of nested contexts in which the event (or
circumstance) is perceived to occur. The terms constant, parameter, and variable,
borrowed from mathematics, seem useful here. Often mathematicians, physicists, and
others will carry out a calculation, saying "c is a constant, p is a parameter, and v is a
variable". What they




Artificial Intelligence:Prospects                                                        640

mean.is that any of them can vary (including the "constant"); however, there is a kind of
hierarchy of variability. In the situation which is being represented b~ the symbols, c
establishes a global condition; p establishes some less global condition which can vary
while c is held fixed; and finally, v can run around while c and p are held fixed. It makes
little sense to think of holding v fixed while c and p vary, for c and p establish the context
in which v has meaning. For instance, think of a dentist who has a list of patients, and for
each patient, a list of teeth. It makes perfect sense (and plenty of money) to hold the
patient fixed and vary his teeth-but it makes no sense at all to hold one tooth fixed and
vary the patient. (Although sometimes it makes good sense to vary the dentist ...)
         We build up our mental representation of a situation layer by layer. The lowest
layer establishes the deepest aspect of the context-sometimes being so low that it cannot
vary at all. For instance, the three-dimensionality of our world is so ingrained that most
of us never would imagine letting it slip mentally. It is a constant constant. Then there are
layers which establish temporarily, though not permanently, fixed aspects of situations,
which could be called background assumptions-things which, in the back of your mind,
you know can vary, but which most of the time you unquestioningly accept as
unchanging aspects. These could still be called "constants". For instance, when you go to
a football game, the rules of the game are constants of that sort. Then there are
"parameters": you think of them as more variable, but you temporarily hold them
constant. At a football game, parameters might include the weather, the opposing team,
and so forth. There could be-and probably are-several layers of parameters. Finally, we
reach the "shakiest" aspects of your mental representation of the situation-the variables.
These are things such as Palindromi's stepping out of bounds, which are mentally "loose"
and which you don't mind letting slip away from their real values, for a short moment.

                            Frames and Nested Contexts
The word frame is in vogue in Al currently, and it could be defined as a computational
instantiation of a context. The term is due to Marvin Minsky, as are many ideas about
frames, though the general concept has been floating around for a good number of years.
In frame language, one could say that mental representations of situations involve frames
nested within each other. Each of the various ingredients of a situation has its own frame.
It is interesting to verbalize explicitly one of my mental images concerning nested
frames. Imagine a large collection of chests of drawers. When you choose a chest, you
have a frame, and the drawer holes are places where "subframes" can be attached. But
subframes are themselves chests of drawers. How can you stick a whole chest of drawers
into the slot for a single drawer in another chest of drawers? Easy: you shrink and distort
the second chest, since, after all, this is all mental, not physical. Now in the outer frame,
there may be several different drawer slots that need to be




Artificial Intelligence:Prospects                                                         641

filled; then you may need to fill slots in some of the inner chests of drawers (or
subframes). This can go on, recursively.
         The vivid surrealistic image of squishing and bending a chest of drawers so that it
can fit into a slot of arbitrary shape is probably quite important, because it hints that your
concepts are squished and bent by the contexts you force them into. Thus, what does your
concept of "person" ', become when the people you are thinking about are football
players? It certainly is a distorted concept, one which is forced on you by the overall
context. You have stuck the "person" frame into a slot in the "football game" frame. The
theory of representing knowledge in frames relies on the idea that the world consists of
quasi-closed subsystems, each of which can serve as a context for others without being
too disrupted, or creating too much disruption, in the process.
         One of the main ideas about frames is that each frame comes with its own set of
expectations. The corresponding image is that each chest. of drawers comes with a built-
in, but loosely bound, drawer in each of its
drawer slots, called a default. If I tell you, "Picture a river bank", you will invoke a visual
image which has various features, most of which you could override if I added extra
phrases such as "in a drought" or "in Brazil" or "without a merry-go-round". The
existence of default values for slots allows the recursive process of filling slots to come to
an end. In effect, you say, "I will fill in the slots myself as far as three layers down;
beyond that I will take the default options." Together with its default expectations, a
frame contains knowledge of its limits of applicability, and heuristics for switching to
other frames in case it has been stretched beyond its limits of tolerance.
         The nested structure of a frame gives you a way of "zooming in" and looking at
small details from as close up as you wish: you just zoom in on the proper subframe, and
then on one of its subframes, etc., until you have the desired amount of detail. It is like
having a road atlas of the USA which has a map of the whole country in the front, with
individual state maps inside, and even maps of cities and some of the larger towns if you
want still more detail. One can imagine an atlas with arbitrary amounts of detail, going
down to single blocks, houses, rooms, etc. It is like looking through a telescope with
lenses of different power; each lens has its own uses. It is important that one can make
use of all the different scales; often detail is irrelevant and even distracting.
         Because arbitrarily different frames can be stuck inside other frames' slots, there
is great potential for conflict or "collision". The nice neat scheme of a single, global set of
layers of "constants", "parameters", and "variables" is an oversimplification. In fact, each
frame will have its own hierarchy of variability, and this is what makes analyzing how we
perceive such a complex event as a football game, with its many subframes,
subsubframes, etc., an incredibly messy operation. How do all these many frames interact
with each other? If there is a conflict where one frame says, "This item is a constant" but
another frame says, "No, it is a variable!", how does it get resolved? These are deep and
difficult problems of frame theory to




Artificial Intelligence:Prospects                                                          642

which I can give no answers. There has as yet been no complete agreement on what a
frame really is, or on how to implement frames in Al programs. I make my own stab at
discussing some of these questions in the following section, where I talk about some
puzzles in visual pattern recognition, which I call "Bongard problems".

                                    Bongard Problems

Bongard problems (BP's) are problems of the general type given by the Russian scientist
M. Bongard in his book Pattern Recognition. A typical BP-number 51 in his collection of
one hundred-is shown in Figure 119.




FIGURE 119. Bongard problem 51. From R1. Bongard, Pattern Recognition (Rochelle
Park, N.,J.: Hayden Book Co., Spartan Books, 1970).]

These fascinating problems are intended for pattern-recognizers, whether human or
machine. (One might also throw in ETI's-extraterrestrial intelligences.) Each problem
consists of twelve boxed figures (henceforth called boxes): six on the left, forming Class
I, and six on the right, forming Class II. The boxes may be indexed this way:

                       I-A    I-B            II-A    II-B
                       I-C    I-D            II-C    II-D
                       I-E    I-F            II-E    II-F

The problem is "How do Class I boxes differ from Class II boxes?"
        A Bongard problem-solving program would have several stages, in which raw
data gradually get converted into descriptions. The early stages are relatively inflexible,
and higher stages become gradually more flexible. The final stages have a property which
I call tentativity, which means simply that the way a picture is represented is always
tentative. Upon the drop of a hat, a high-level description can be restructured, using all
the devices of the




Artificial Intelligence:Prospects                                                      643

later stages. The ideas presented below also have a tentative quality to them. I will try to
convey overall ideas first, glossing over significant difficulties. Then I will go back and
try to explain subtleties and tricks and so forth. So your notion of how it all works may
also undergo some revisions as you read. But that is in the spirit of the discussion.

                    Preprocessing Selects a Mini-vocabulary
Suppose, then, that we have some Bongard problem which we want to solve. The
problem is presented to a TV camera and the raw data are read in. Then the raw data are
preprocessed. This means that some salient features are detected. The names of these
features constitute a "mini-vocabulary" for the problem; they are drawn from a general
"salient-feature vocabulary". Some typical terms of the salient-feature vocabulary are:

       line segment, curve, horizontal, vertical, black, white, big, small, pointy, round ...

In a second stage of preprocessing, some knowledge about elementary shapes is used;
and if any are found, their names are also made available. Thus, terms such as

       triangle, circle, square, indentation, protrusion, right angle, vertex, cusp, arrow ...

may be selected. This is roughly the point at which the conscious and the unconscious
meet, in humans. This discussion is primarily concerned with describing what happens
from here on out.

                                High-Level Descriptions
Now that the picture is "understood", to some extent, in terms of familiar concepts, some
looking around is done. Tentative descriptions are made for one or a few of the twelve
boxes. They will typically use simple descriptors such as

        above, below, to the right of, to the left of, inside, outside of, close to, far from,
parallel to, perpendicular to, in a row, scattered, evenly spaced, irregularly spaced, etc.

Also, definite and indefinite numerical descriptors can be used:

       1, 2, 3, 4, 5, ... many, few, etc.

More complicated descriptors may be built up, such as

       further to the right of, less close to, almost parallel to, etc.




Artificial Intelligence:Prospects                                                          644

       FIGURE 120. Bongard problem 47. [From M. Bongard, Pattern Recognition.]

Thus, a typical box-say I-F of BP 47 (Fig. 120)-could be variously described as having:

                                             three shapes
                                                  or
                                         three white shapes
                                                  or
                                        a circle on the right
                                                  or
                                      two triangles and a circle
                                                  or
                                  two upwards-pointing triangles
                                                  or
                               one large shape and two small shapes
                                                  or
                          one curved shape and two straight-edged shapes
                                                  or
                  a circle with the same kind of shape on the inside and outside.

Each of these descriptions sees the box through a "filter". Out of context, any of them
might be a useful description. As it turns out, though, all of them are "wrong", in the
context of the particular Bongard problem they are part of. In other words, if you knew
the distinction between Classes I and II in BP 47, and were given one of the preceding
lines as a description of an unseen drawing, that information would not allow you to tell
to which Class the drawing belonged. The essential feature of this box, in context, is that
it includes

                               a circle containing a triangle.

Note that someone who heard such a description would not be able to reconstruct the
original drawing, but would be able to recognize drawings




Artificial Intelligence:Prospects                                                      645

FIGURE 121. Bongard problem 91. [From M. Bongard, Pattern Recognition.]

which have this property. It is a little like musical style: you may be an infallible
recognizer of Mozart, but at the same time unable to write anything which would fool
anybody into thinking it was by Mozart.
        Now consider box I-D of BP 91 (Fig. 121). An overloaded but "right" description
in the context of BP 91 is

                          a circle with three rectangular intrusions.

Notice the sophistication of such a description, in which the word "with" functions as a
disclaimer, implying that the "circle" is not really a circle: it is almost a circle, except that
. . . Furthermore, the intrusions are not full rectangles. There is a lot of "play" in the way
we use language to describe
         things. Clearly, a lot of information has been thrown away, and even more could
be thrown away. A priori, it is very hard to know what it would be smart to throw away
and what to keep. So some sort of method for an intelligent compromise has to be
encoded, via heuristics. Of course, there is always recourse to lower levels of description
(i.e., less chunked descriptions) if discarded information has to be retrieved, just as
people can constantly look at the puzzle for help in restructuring their ideas about it. The
trick, then, is to devise explicit rules that say how to

   make tentative descriptions for each box;
   compare them with tentative descriptions for other boxes of either Class;
   restructure the descriptions, by
       (i) adding information,
       (ii) discarding information,
   or (iii) viewing the same information from another angle; iterate this process until
   finding out what makes the two Classes differ.




Artificial Intelligence:Prospects                                                           646

   Templates and Sameness-Detectors
   One good strategy would be to try to make descriptions structurally similar to
   each other, to the extent this is possible. Any structure they have in common will
   make comparing them that much easier. Two important elements of this theory
   deal with this strategy. One is the idea of "description-schemas" or templates; the
   other is the idea of Sam-a "sameness detector".
       First Sam. Sam is a special agent present on all levels of the program.
   (Actually there may be different kinds of Sams on different levels.) Sam
   constantly runs around within individual descriptions and within different
   descriptions, looking for descriptors or other things which are repeated. When
   some sameness is found, various restructuring operations can be triggered, either
   on the single-description level or on the level of several descriptions at once.
   Now templates. The first thing that happens after preprocessing is an attempt to
   manufacture a template, or description-schema-a un form format for the
   descriptions of all the boxes in a problem. The idea is that a description can often
   be broken up in a natural way into subdescriptions, and those in turn into subs
   ubdescriptions, if need be. The bottom is hit when you come to primitive concepts
   which belong to the level of the preprocessor. Now it is important to choose the
   way of breaking descriptions into parts so as to reflect commonality among all the
   boxes; otherwise you are introducing a superfluous and meaningless kind of
   "pseudo-order" into the world.
       On the basis of what information is a template built? It is best to look at an
   example. Take BP 49 (Fig. 122). Preprocessing yields the information that each
   box consists of several little o's, and one large closed curve. This is a valuable
   observation, and deserves to be incorporated in the template. Thus a first stab at a
   template would be:

   large closed curve:-----
   small o's:-----

   FIGURE 122. Bongard problem 49. [From M. Bongard, Pattern Recognition.]




Artificial Intelligence:Prospects                                                     647

   It is very simple: the description-template has two explicit slots where
   subdescriptions are to be attached.

                                 A Heterarchical Program
   Sow an interesting thing happens, triggered by the term "closed curve". one of the
   most important modules in the program is a kind of semantic net--the concept
   network-in which all the known nouns, adjectives, etc., are linked in ways which
   indicate their interrelations. For instance, "closed curve" is strongly linked with
   the terms "interior" and "exterior". The concept net is just brimming with
   information about relations between terms, such as what is the opposite of what,
   what is similar to what, what often occurs with what, and so on. A little portion of
   a concept network, to be explained shortly, is shown in Figure 123. But let us
   follow what happens now, in the solution of problem 49. The concepts "interior"
   and "exterior" are activated by their proximity in the net to "closed curve". This
   suggests to the template-builder that it might be a good idea to make distinct slots
   for the interior and exterior of the curve. Thus, in the spirit of tentativity, the
   template is tentatively restructured to be this:

       large closed curve: ----
       little o's in interior: ----
       little o's in exterior:----

   Now when subdescriptions are sought, the terms "interior" and "exterior" will
   cause procedures to inspect those specific regions of the box. What is found in BP
   49, box I-A is this:

       large closed curve: circle
       little o's in interior: three
       little o's in exterior: three

   And a description of box II-A of the same BP might be

       large closed curve: cigar
       little o's in interior: three
       little o's in exterior: three

       Now Sam, constantly active in parallel with other operations, spots the
   recurrence of the concept "three" in all the slots dealing with o's, and this is strong
   reason to undertake a second template-restructuring operation. Notice that the first
   was suggested by the concept net, the second by Sam. Now our template for
   problem 49 becomes:

       large closed curve:----
       three little o's in interior: -----
       three little o's in exterior:-----


Artificial Intelligence:Prospects                                                        648

   FIGURE 123. A small portion of a concept network for a program to solve
   Bongard Problems. "Nodes" are joined by "links", which in turn can be linked. By
   considering a link as a verb and the nodes it joins as subject and object, you can
   pull out some English sentences from this diagram.



Artificial Intelligence:Prospects                                                   649

   Now that "three" has risen" one level of generality-namely, into the template-it
   becomes worthwhile to explore its neighbors in the concept network. One of them
   is "triangle", which suggests that triangles of o's may be important. As it happens,
   this leads down a blind alley-but how could you know in advances It is a typical
   blind alley that a human would explore, so it is good if our program finds it too!
   For box II-E, a description such as the following might get generated:

   large closed curve: circle
   three little o's in interior: equilateral triangle
   three little o's in exterior: equilateral triangle

   Of course an enormous amount of information has been thrown away concerning
   the sizes, positions, and orientations of these triangles, and many other things as
   well. But that is the whole point of making descriptions instead of just using the
   raw data! It is the same idea as funneling, which we discussed in Chapter XI.

                                 The Concept Network
   We need not run through the entire solution of problem 49; this suffices to show
   the constant back-and-forth interaction of individual descriptions, templates, the
   sameness-detector Sam, and the concept network. We should now look a little
   more in detail at the concept network and its function. A simplified portion shown
   in the figure codes the following ideas:

       "High" and "low" are opposites.
       "Up" and "down" are opposites.
       "High" and "up" are similar.
       "Low" and "down" are similar.
       "Right" and "left" are opposites.
       The "right-left" distinction is similar to the "high-low" distinction.
       "Opposite" and "similar" are opposites.

   Note how everything in the net-both nodes and links-can be talked about. In that
   sense nothing in the net is on a higher level than anything else. Another portion of
   the net is shown; it codes for the ideas that

       A square is a polygon.
       A triangle is a polygon.
       A polygon is a closed curve.




Artificial Intelligence:Prospects                                                     650

       The difference between a triangle and a square is that one has 3 sides and the
           other has 4.
       4 is similar to 3.
       A circle is a closed curve.
       A closed curve has an interior and an exterior.
       "Interior" and "exterior" are opposites.

   The network of concepts is necessarily very vast. It seems to store knowledge only
   statically, or declaratively, but that is only half the story. Actually, its knowledge
   borders on being procedural as well, by the fact that the proximities in the net act
   as guides, or "programs", telling the main program how to develop its
   understanding of the drawings in the boxes.
   For instance, some early hunch may turn out to be wrong and yet have the germ of
   the right answer in it. In BP 33 (Fig. 124), one might at first




   FIGURE 124. Bongard problem 33. [From M. Bongard, Pattern Recognition.]


   jump to the idea that Class I boxes contain "pointy" shapes, Class II boxes contain
   "smooth" ones. But on closer inspection, this is wrong. Nevertheless, there is a
   worthwhile insight here, and one can try to push it further, by sliding around in the
   network of concepts beginning at "pointy". It is close to the concept "acute",
   which is precisely the distinguishing feature of Class I. Thus one of the main
   functions of the concept network is to allow early wrong ideas to be modified
   slightly, to slip into variations which may be correct.

                              Slippage and Tentativity
   Related to this notion of slipping between closely related terms is the notion of
   seeing a given object as a variation on another object. An excellent example has
   been mentioned already-that of the "circle with three indentations", where in fact
   there is no circle at all. One has to be able to bend concepts, when it is appropriate.
   Nothing should be absolutely rigid. On


Artificial Intelligence:Prospects                                                        651

   the other hand, things shouldn't be so wishy-washy that nothing has any meaning
   at all, either. The trick is to know when and how to slip one concept into another.
        An extremely interesting set of examples where slipping from one description
   to another is the crux of the matter is given in Bongard problems 85-87 (Fig. 125).
   BP 85 is rather trivial. Let us assume that our program identifies "line segment" in
   its preprocessing stage. It is relatively simple for it then to count line segments and
   arrive at the difference

   FIGURE 125.Bongard problems 85-87. [From M. Bongard, Pattern Recognition.]




Artificial Intelligence:Prospects                                                        652

   between Class I and Class II in BP 85. Now it goes on to BP 86. A general
   heuristic which it uses is to try out recent ideas which have worked. Successful
   repetition of recent methods is very common in the real world, and Bongard does
   not try to outwit this kind of heuristic in his collection-in fact, he reinforces it,
   fortunately. So we plunge right into problem 86 with two ideas ("count" and "line
   segment") fused into one: "count line segments". But as it happens, the trick of BP
   86 is to count line trains rather than line segments, where "line train" means an
   end-to-end concatenation of (one or more) line segments. One way the program
   might figure this out is if the concepts "line train" and "line segment" are both
   known, and are close in the concept network. Another way is if it can invent the
   concept of "line train"-a tricky proposition, to say the least.
       Then comes BP 87, in which the notion of "line segment" is further played
   with. When is a line segment three line segments? (See box II-A.) The program
   must be sufficiently flexible that it can go back and forth between such different
   representations for a given part of a drawing. It is wise to store old representations,
   rather than forgetting them and perhaps having to reconstruct them, for there is no
   guarantee that a newer representation is better than an old one. Thus, along with
   each old representation should be stored some of the reasons for liking it and
   disliking it. (This begins to sound rather complex, doesn't it?)

                                    Meta- Descriptions
   Now we come to another vital part of the recognition process, and that has to do
   with levels of abstraction and meta-descriptions. For this let us consider BP 91
   (Fig. 121) again. What kind of template could be constructed here? There is such
   an amount of variety that it is hard to know where to begin. But this is in itself a
   clue! The clue says, namely, that the class distinction very likely exists on a higher
   level of abstraction than that of geometrical description. This observation clues the
   program that it should construct descriptions of descriptions-that is, meta-
   descriptions. Perhaps on this second level some common feature will emerge; and
   if we are lucky, we will discover enough commonality to guide us towards the
   formulation of a template for the meta-descriptions! So we plunge ahead without a
   template, and manufacture descriptions for various boxes; then, once these
   descriptions have been made, we describe them. What kinds of slot will our
   template for meta-descriptions have? Perhaps these, among others:

   concepts used: ----
   recurring concepts-----:
   names of slots: -----
   filters used:----

   There are many other kinds of slots which might be needed in metadescriptions,
   but this is a sample. Now suppose we have described box I-E of BP 91. Its
   (template-less) description might look like this:




Artificial Intelligence:Prospects                                                        653

       horizontal line segment
       vertical line segment mounted on the horizontal line segment
       vertical line segment mounted on the horizontal line segment
       vertical line segment mounted on the horizontal line segment

   Of course much information has been thrown out: the fact that the three vertical
   lines are of the same length, are spaced equidistantly, etc. But it is plausible that
   the above description would be made. So the meta description might look like this:

       concepts used: vertical-horizontal, line segment, mounted on
       repetitions in description: 3 copies of "vertical line segment mounted on the
       horizontal line segment"
       names of slots-----
       filters used:-----

   Not all slots of the meta-description need be filled in; information can be thrown
   away on this level as well as on the Just- plain-description" level.
   ‘Now if we were to make a description for any of the other boxes of Class I, and
   then a meta-description of it, we would wind up filling the slot "repetitions in
   description" each time with the phrase "3 copies of ..." The sameness-detector
   would notice this, and pick up three-ness as a salient feature, on quite a high level
   of abstraction, of the boxes of Class I. Similarly, four-ness would be recognized,
   via the method of metadescriptions, as the mark of Class II.

                               Flexibility Is Important
   Now you might object that in this case, resorting to the method of
   metadescriptions is like shooting a fly with an elephant gun, for the three-ness
   versus four-ness might as easily have shown up on the lower level if we had
   constructed our descriptions slightly differently. Yes, true-but it is important to
   have the possibility of solving these problems by different routes. There should be
   a large amount of flexibility in the program; it should not be doomed if,
   malaphorically speaking, it "barks up the wrong alley" for a while. (The amusing
   term "malaphor" was coined by the newspaper columnist Lawrence Harrison; it
   means a cross between a malapropism and a metaphor. It is a good example of
   "recombinant ideas".) In any case, I wanted to illustrate the general principle that
   says: When it is hard to build a template because the preprocessor finds too much
   diversity, that should serve as a clue that concepts on a higher level of abstraction
   are involved than the preprocessor knows about.

                               Focusing and Filtering
   Now let us deal with another question: ways to throw information out. This
   involves two related notions, which I call "focusing" and "filtering". Focus-




Artificial Intelligence:Prospects                                                      654

   FIGURE 126. Bongard problem 55. [From M. Bongard, Pattern Recognition.]




   FIGURE 127. Bongard problem 22. [From M. Bongard, Pattern Recognition.]

   ing involves making a description whose focus is some part of the drawing in the
   box, to the exclusion of everything else. Filtering involves making a description
   which concentrates on some particular way of viewing the contents of the box, and
   deliberately ignores all other aspects. Thus they are complementary: focusing has
   to do with objects (roughly, nouns), and filtering has to do with concepts (roughly,
   adjectives). For an example of focusing, let's look at BP 55 (Fig. 126). Here, we
   focus on the indentation and the little circle next to it, to the exclusion of the
   everything else in the box. BP 22 (Fig. 127) presents an example of filtering. Here,
   we must filter out every concept but that of size. A combination of focusing and
   filtering is required to solve problem BP 58 (Fig. 128).
   One of the most important ways to get ideas for focusing and filtering is by
   another sort of "focusing": namely, by inspection of a single particularly simple
   box-say one with as few objects in it as possible. It can be




Artificial Intelligence:Prospects                                                     655

   FIGURE 128. Bongard problem 58. [From M. Bongard, Pattern Recognition.]




   FIGURE 129. Bongard problem 61. [From M. Bongard, Pattern Recognition.]

   extremely helpful to compare the starkest boxes from the two Classes. But how
   can you tell which boxes are stark until you have descriptions for them? Well, one
   way of detecting starkness is to look for a box with a minimum of the features
   provided by the preprocessor. This can be done very early, for it does not require a
   pre-existing template; in fact, this can be one useful way of discovering features to
   build into a template. BP 61 (Fig. 129) is an example where that technique might
   quickly lead to a solution.

                 Science and the World of Bongard Problems
   One can think of the Bongard-problem world as a tiny place where "science" is
   done-that is, where the purpose is to discern patterns in the world. As patterns are
   sought, templates are made, unmade, and remade;




Artificial Intelligence:Prospects                                                      656

   FIGURE 130. Bongard problems 70-71. [From M. Bongard, Pattern
   Recognition.]

   slots are shifted from one level of generality to another: filtering and focusing are
   done; and so on. There are discoveries on all levels of complexity. The Kuhnian
   theory that certain rare events called "paradigm shifts" mark the distinction
   between "normal" science and "conceptual revolutions" does not seem to work,
   for we can see paradigm shifts happening all throughout the system, all the time.
   The fluidity of descriptions ensures that paradigm shifts will take place on all
   scales.
       Of course, some discoveries are more "revolutionary" than others, because
   they have wider effects. For instance, one can make the discovery that problems
   70 and 71 (Fig. 130) are "the same problem", when looked at on a sufficiently
   abstract level. The key observation is that both involve depth-2 versus depth-1
   nesting. This is a new level of discovery that can he made about Bongard
   problems. There is an even higher level, concerning the collection as a whole. If
   someone has never seen the collection, it can be a good puzzle just to figure out
   what it is. To figure it out is a revolutionary insight, but it must be pointed out that
   the mechanisms of thought which allow such a discovery to be made are no
   different from those which operate in the solution of a single Bongard problem.




Artificial Intelligence:Prospects                                                         657

By the same token, real science does not divide up into "normal" periods versus
"conceptual. revolutions"; rather, paradigm shifts pervade-there are just bigger and
smaller ones, paradigm shifts on different levels. The recursive plots of INT and
Gplot (Figs. 32 and 34) provide a geometric model for this idea: they have the same
structure full of discontinuous jumps on every level, not just the top level-only the
lower the level, the smaller the jumps

                   Connections to Other Types of Thought
To set this entire program somewhat in context, let me suggest two ways in which it is
related to other aspects of cognition. Not only does it depend on other aspects of
cognition, but also they in turn depend on it. First let me comment on how it depends on
other aspects of cognition. The intuition which is required for knowing when it makes
sense to blur distinctions, to try redescriptions, to backtrack, to shift levels, and so forth,
is something which probably comes only with much experience in thought in general.
Thus it would be very hard to define heuristics for these crucial aspects of the program.
Sometimes one's experience with real objects in the world has a subtle effect on how one
describes or redescribes boxes. For instance, who can say how much one's familiarity
with living trees helps one to solve BP 70\% It is very doubtful that in humans, the
subnetwork of concepts relevant to these puzzles can be easily separated out from the
whole network. Rather, it is much more likely that one's intuitions gained from seeing
and handling real objects-combs, trains, strings, blocks, letters, rubber bands, etc., etc.-
play an invisible but significant guiding role in the solution of these puzzles.
        Conversely, it is certain that understanding real-world situations heavily depends
on visual imagery and spatial intuition, so that having a powerful and flexible way of
representing patterns such as these Bongard patterns can only contribute to the general
efficiency of thought processes.
It seems to me that Bongard's problems were worked out with great care, and that they
have a quality of universality to them, in the sense that each one has a unique correct
answer. Of course one could argue with this and say that what we consider "correct"
depends in some deep way on our being human, and some creatures from some other star
system might disagree entirely. Not having any concrete evidence either way, I still have
a certain faith that Bongard problems depend on a sense of simplicity which is not just
limited to earthbound human beings. My earlier comments about the probable importance
of being acquainted with such surely earth-limited objects as combs, trains, rubber bands,
and so on, are not in conflict with the idea that our notion of simplicity is universal, for
what matters is not any of these individual objects, but the fact that taken together they
span a wide space. And it seems likely that any other civilization would have as vast a
repertoire of artifacts and natural objects and varieties of experience on which to draw as
we do. So I believe that the skill of solving Bongard




Artificial Intelligence:Prospects                                                          658

problems lies very close to the core of "pure" intelligence, if there is such a thing.
Therefore it is a good place to begin if one wants to investigate the ability to discover
"intrinsic meaning" in patterns or messages. Unfortunately we have reproduced only a
small selection of his stimulating collection. I hope that many readers will acquaint
themselves with the entire collection, to be found in his book (see Bibliography).
        Some of the problems of visual pattern recognition which we human beings seem
to have completely "flattened" into our unconscious are quite amazing. They include:

   recognition of faces (invariance of faces under age change, expression change,
      lighting change, distance change, angle change, etc.)
   recognition of hiking trails in forests and mountains-somehow this has always
      impressed me as one of our most subtle acts of pattern recognition-and yet
      animals can do it, too
   reading text without hesitation in hundreds if not thousands of different typefaces

             Message-Passing Languages, Frames, and Symbols
One way that has been suggested for handling the complexities of pattern recognition and
other challenges to Al programs is the so-called "actor" formalism of Carl Hewitt (similar
to the language "Smailtalk", developed by Alan Kay and others), in which a program is
written as a collection of interacting actors, which can pass elaborate messages back and
forth among themselves. In a way, this resembles a heterarchical collection of procedures
which can call each other. The major difference is that where procedures usually only
pass a rather small number of arguments back and forth, the messages exchanged by
actors can be arbitrarily long and complex.
         Actors with the ability to exchange messages become somewhat autonomous
agents-in fact, even like autonomous computers, with messages being somewhat like
programs. Each actor can have its own idiosyncratic way of interpreting any given
message; thus a message's meaning will depend on the actor it is intercepted by. This
comes about by the actor having within it a piece of program which interprets messages;
so there may be as many interpreters as there are actors. Of course, there may be many
actors with identical interpreters; in fact, this could be a great advantage, just as it is
extremely important in the cell to have a multitude of identical ribosomes floating
throughout the cytoplasm, all of which will interpret a message-in this case, messenger
RNA-in one and the same way.
It is interesting to think how one might merge the frame-notion with the actor-notion. Let
us call a frame with the capability of generating and interpreting complex messages a
symbol:

                                    frame + actor = symbol




Artificial Intelligence:Prospects                                                      659

We now have reached the point where we are talking about ways or implementing those
elusive active symbols of Chapters XI and XII; henceforth in this Chapter, "symbol" will
have that meaning. By the way, don't feel dumb if you don't immediately see just how
this synthesis is to be made. It is not clear, though it is certainly one of the most
fascinating directions to go in AI. Furthermore, it is quite certain that even the best
synthesis of these notions will turn out to have much less power than the actual symbols
of human minds. In that sense, calling these frame-actor syntheses "symbols" is
premature, but it is an optimistic way of looking at things.
        Let us return to some issues connected with message passing. Should each
message be directed specifically at a target symbol, or should it be thrown out into the
grand void, much as mRNA is thrown out into the cytoplasm, to seek its ribosome? If
messages have destinations, then each symbol must have an address, and messages for it
should always be sent to that address. On the other hand, there could be one central
receiving dock for messages, where a message would simply sit until it got picked up by
some symbol that wanted it. This is a counterpart to General Delivery. Probably the best
solution is to allow both types of message to exist; also to have provisions for different
classes of urgency-special delivery, first class, second class, and so on. The whole postal
system provides a rich source of ideas for message-passing languages, including such
curios as selfaddressed stamped envelopes (messages whose senders want answers
quickly), parcel post (extremely long messages which can be sent some very slow way),
and more. The telephone system will give you more inspiration when you run out of
postal-system ideas.

                                    Enzymes and AI
Another rich source of ideas for message passing-indeed, for information processing in
general-is, of course, the cell. Some objects in the cell are quite comparable to actors-in
particular, enzymes. Each enzyme's active site acts as a filter which only recognizes
certain kinds of substrates (messages). Thus an enzyme has an "address", in effect. The
enzyme is "programmed" (by virtue of its tertiary structure) to carry out certain
operations upon that "message", and then to release it to the world again. Now in this
way, when a message is passed from enzyme to enzyme along a chemical pathway, a lot
can be accomplished. We have already described the elaborate kinds of feedback
mechanisms which can take place in cells (either by inhibition or repression). These kinds
of mechanisms show that complicated control of processes can arise through the kind of
message passing that exists in the cell.
        One of the most striking things about enzymes is how they sit around idly,
waiting to be triggered by an incoming substrate. Then, when the substrate arrives,
suddenly the enzyme springs into action, like a Venus's flytrap. This kind of "hair-
trigger" program has been used in Al, and goes by the name of demon. The important
thing here is the idea of having many different "species" of triggerable subroutines just
lying around waiting to




Artificial Intelligence:Prospects                                                      660

be triggered. In cells, all the complex molecules and organelles are built up, simple step
by simple step. Some of these new structures are often enzymes themselves, and they
participate in the building of new enzymes, which in turn participate in the building of yet
other types of enzyme, etc. Such recursive cascades of enzymes can have drastic effects
on what a cell is doing. One would like to see the same kind of simple step-by-step
assembly process imported into AI, in the construction of useful subprograms. For
instance, repetition has a way of burning new circuits into our mental hardware, so that
oft-repeated pieces of behavior become encoded below the conscious level. It would be
extremely useful if there were an analogous way of synthesizing efficient pieces of code
which can carry out the same sequence of operations as something which has been
learned on a higher level of "consciousness". Enzyme cascades may suggest a model for
how this could be done. (The program called "HACKER", written by Gerald Sussman,
synthesizes and debugs small subroutines in a way not too much unlike that of enzyme
cascades.)
        The sameness-detectors in the Bongard problem-solver (Sams) could be
implemented as enzyme-like subprograms. Like an enzyme, a Sam would meander about
somewhat at random, bumping into small data structures here and there. Upon filling its
two "active sites" with identical data structures, the Sam would emit a message to other
parts (actors) of the program. As long as programs are serial, it would not make much
sense to have several copies of a Sam, but in a truly parallel computer, regulating the
number of copies of a subprogram would be a way of regulating the expected waiting-
time before an operation gets done, just as regulating the number of copies of an enzyme
in a cell regulates how fast that function gets performed. And if new Sams could be
synthesized, that would be comparable to the seepage of pattern detection into lower
levels of our minds.

                                    Fission and Fusion
Two interesting and complementary ideas concerning the interaction of symbols are
"fission" and "fusion". Fission is the gradual divergence of a new symbol from its parent
symbol (that is, from the symbol which served as a template off of which it was copied).
Fusion is what happens when two (or more) originally unrelated symbols participate in a
"joint activation", passing messages so tightly back and forth that they get bound together
and the combination can thereafter be addressed as if it were a single symbol. Fission is a
more or less inevitable process, since once a new symbol has been "rubbed off" of an old
one, it becomes autonomous, and its interactions with the outside world get reflected in
its private internal structure; so what started out as a perfect copy will soon become
imperfect, and then slowly will become less and less like the symbol off of which it was
"rubbed". Fusion is a subtler thing. When do two concepts really become 'one? Is there
some precise instant when a fusion takes place?




Artificial Intelligence:Prospects                                                       661

        This notion of joint activations opens up a Pandora's box of questions. For
instance, how much coo we hear "dough" and "nut" when we say "doughnut"? Does a
German who thinks of gloves ("Handschuhe") hear "hand-shoes" or not? How about
Chinese people, whose word "dong-xi" ("East-West") means "thing"? It is a matter of
some political concern, too, since some people claim that words like "chairman" are
heavily charged with undertones of the male gender. The degree to which the parts
resonate inside the whole probably varies from person to person and according to
circumstances.
        The real problem with this notion of "fusion" of symbols is that it is very hard to
imagine general algorithms which will create meaningful new symbols from colliding
symbols. It is like two strands of DNA which come together. How do you take parts from
each and recombine them into a meaningful and viable new strand of DNA which codes
for an individual of the same species? Or a new kind of species? The chance is
infinitesimal that a random combination of pieces of DNA will code for anything that
will survive-something like the chance that a random combination of words from two
books will make another book. The chance that recombinant DNA will make sense on
any level but the lowest is tiny, precisely because there are so many levels of meaning in
DNA. And the same goes for "recombinant symbols".

                           Epigenesis of the Crab Canon
I think of my Dialogue Crab Canon as a prototype example where two ideas collided in
my mind, connected in a new way, and suddenly a new kind of verbal structure came
alive in my mind. Of course I can still think about musical crab canons and verbal
dialogues separately-they can still be activated independently of each other; but the fused
symbol for crab canonical dialogues has its own characteristic modes of activation, too.
To illustrate this notion of fusion or "symbolic recombination" in some detail, then, I
would like to use the development of my Crab Canon as a case study, because, of course,
it is very familiar to me, and also because it is interesting, yet typical of how far a single
idea can be pushed. I will recount it in stages named after those of meiosis, which is the
name for cell division in which "crossing-over", or genetic recombination, takes place-the
source of diversity in evolution.
         PROPHASE: I began with a rather simple idea-that a piece of music, say a canon,
could be imitated verbally. This came from the observation that, through a shared abstract
form, a piece of text and a piece of music may be connected. The next step involved
trying to realize some of the potential of this vague hunch; here, I hit upon the idea that
"voices" in canons can be mapped onto "characters" in dialogues-still a rather obvious
idea.
         Then I focused down onto specific kinds of canons, and remembered that there
was a crab canon in the Musical Offering. At that time, I had just




Artificial Intelligence:Prospects                                                         662

begun writing Dialogues, and there were only two characters: Achilles and the Tortoise.
Since the Bach crab canon has two voices, this mapped perfectly: Achilles should be one
voice, the Tortoise the other, with the one doing forwards what the other does backwards.
But here I was faced with a problem: on what level should the reversal take place? The
letter level? The word level? The sentence level? After some thought, I concluded that
the "dramatic line" level would be most appropriate.
         Now that the "skeleton" of the Bach crab canon had been transplanted, at least in
plan, into a verbal form, there was just one problem. When the two voices crossed in the
middle, there would be a short period of extreme repetition: an ugly blemish. What to do
about it? Here, a strange thing happened, a kind of level-crossing typical of creative acts:
the word "crab" in "crab canon" flashed into my mind, undoubtedly because of some
abstract shared quality with the notion of "tortoise"-and immediately I realized that at the
dead center, I could block the repetitive effect, by inserting one special line, said by a
new character: a Crab! This is how, in the "prophase" of the Crab Canon, the Crab was
conceived: at the crossing over of Achilles and the Tortoise. (See Fig. 131.)




FIGURE 131. A schematic diagram of the Dialogue Crab Canon.

         METAPHASE: This was the skeleton of my Crab Canon. I then entered the
second stage-the "metaphase"-in which I had to fill in the flesh, which was of course an
arduous task. I made a lot of stabs at it, getting used to the way in which pairs of
successive lines had to make sense when read from either direction, and experimenting
around to see what kinds of dual meanings would help me in writing such a form (e.g.,
"Not at all"). There were two early versions both of which were interesting, but weak. I
abandoned work on the book for over a year, and when I returned to the Crab Canon, I
had a few new ideas. One of them was to mention a Bach canon inside it. At first my plan
was to mention the "Canon per augmentationem, contrario motu", from the Musical
Offering (Sloth Canon, as I call it). But that started to seem a little silly, so reluctantly I
decided that inside my Crab Canon, I could talk about Bach's own Crab Canon instead.
Actually, this was a crucial turning point, but I didn't know it then.
         Now if one character was going to mention a Bach piece, wouldn't it be awkward
for the other to say exactly the same thing in the corresponding place? Well, Escher was
playing a similar role to Bach in my thoughts and my book, so wasn't there some way of
just slightly modifying the line so that it would refer to Escher? After all, in the strict art
of canons, note-perfect imitation is occasionally foregone for the sake of elegance or
beauty. And




Artificial Intelligence:Prospects                                                          663

no sooner did that idea occur to me than the picture Day and Night (Fig. 49) popped into
my mind. "Of course!" I thought, "It is a sort of pictorial crab canon, with essentially two
complementary voices carrying the same theme both leftwards and rightwards, and
harmonizing with each other!" Here again was the notion of a single "conceptual
skeleton" being instantiated in two different media-in this case, music and art. So I let the
Tortoise talk about Bach, and Achilles talk about Escher, in parallel language; certainly
this slight departure from strict imitation retained the spirit of crab cano.is.
         At this point, I began realizing that something marvelous was happening namely,
the Dialogue was becoming self-referential, without my even having intended it! What's
more, it was an indirect self-reference, in that the characters did not talk directly about
the Dialogue they were in, but rather about structures which were isomorphic to it (on a
certain plane of abstraction). To put it in the terms I have been using, my Dialogue now
shared a "conceptual skeleton" with Gödel’s G, and could therefore be mapped onto G in
somewhat the way that the Central Dogma was, to create in this case a "Central
Crabmap". This was most exciting to me, since out of nowhere had come an esthetically
pleasing unity of Gödel, Escher, and Bach.
         ANAPHASE: The next step was quite startling. I had had Caroline MacGillavry's
monograph on Escher's tessellations for years, but one day, as I flipped through it, my
eye was riveted to Plate 23 (Fig. 42), for I saw it in a way I had never seen it before: here
was a genuine crab canon-crab-like in both form and content! Escher himself had given
the picture no title, and since he had drawn similar tessellations using many other animal
forms, it is probable that this coincidence of form and content was just something which I
had noticed. But fortuitous or not, this untitled plate was a miniature version of one main
idea of my book: to unite form and content. So with delight I christened it Crab Canon,
substituted it for Day and Night, and modified Achilles' and the Tortoise's remarks
accordingly.
         Yet this was not all. Having become infatuated with molecular biology, one day I
was perusing Watson's book in the bookstore, and in the index saw the word
"palindrome". When I looked it up, I found a magical thing: crab-canonical structures in
DNA. Soon the Crab's comments had been suitably modified to include a short remark to
the effect that he owed his predilection for confusing retrograde and forward motion to
his genes.
         TELOPHASE: The last step came months later, when, as I was talking about the
picture of the crab-canonical section of DNA (Fig. 43), 1 saw that the 'A', 'T', 'C' of
Adenine, Thymine, Cytosine coincided- mirabile dictu-with the 'A', 'T', 'C' of Achilles,
Tortoise, Crab; moreover, just as Adenine and Thymine are paired in DNA, so are
Achilles and the Tortoise paired in the Dialogue. I thought for a moment and, in another
of those level-crossings, saw that 'G', the letter paired with 'C' in DNA, could stand for
"Gene". Once again, I jumped back to the Dialogue, did a little surgery on the Crab's
speech to reflect this new discovery, and now I had a mapping between the DNA's
structure, and the Dialogue's structure. In that sense, the DNA could be said to be a
genotype coding for a phenotype: the




Artificial Intelligence:Prospects                                                        664

Structure of the Dialogue. This final touch dramatically heightened the self-reference,
and gave the Dialogue a density of meaning which I had never anticipated.

               Conceptual Skeletons and Conceptual Mapping
That more or less summarizes the epigenesis of the Crab Canon. The whole process can
be seen as a succession of mappings of ideas onto each other, at varying levels of
abstraction. This is what I call conceptual mapping, and the abstract structures which
connect up two different ideas are conceptual skeletons. Thus, one conceptual skeleton is
that of the abstract notion of a crab canon:

                  a structure having two parts which do the same thing,
                            only moving in opposite directions.

       This is a concrete geometrical image which can be manipulated by the mind
almost as a Bongard pattern. In fact, when I think of the Crab Canon today, I visualize it
as two strands which cross in the middle, where they are joined by a "knot" (the Crab's
speech). This is such a vividly pictorial image that it instantaneously maps, in my mind,
onto a picture of two homologous chromosomes joined by a centromere in their middle,
which is an image drawn directly from meiosis, as shown in Figure 132.




                                       FIGURE 132.

In fact, this very image is what inspired me to cast the description of the Crab Canon's
evolution in terms of meiosis-which is itself, of course, vet another example of
conceptual mapping.

                                    Recombinant Ideas
There are a variety of techniques of fusion of two symbols. One involves lining the two
ideas up next to each other (as if ideas were linear!), then judiciously choosing pieces
from each one, and recombining them in a new symbol. This strongly recalls genetic
recombination. Well, what do chromosomes exchange, and how do they do it? They
exchange genes. What in a symbol is comparable to a gene? If symbols have frame-like
slots, then slots, perhaps. But which slots to exchange, and why? Here is where the
crabcanonical fusion may offer some ideas. Mapping the notion of "musical crab canon"
onto that of "dialogue" involved several auxiliary mappings; in




Artificial Intelligence:Prospects                                                     665

fact it induced them. That is, once it had been decided that these two notions ,ere to be
fused, it became a matter of looking at them on a level where analogous parts emerged
into view, then going ahead and mapping the parts onto each other, and so on,
recursively, to any level that was found desirable. Here, for instance, "voice" and
"character" emerged as corresponding slots when "crab canon" and "dialogue" were
viewed abstractly. Where did these abstract views come from, though? This is at the crux
of the mapping-problem-where do abstract views come from? How do you make abstract
views of specific notions?

                        Abstractions, Skeletons, Analogies
A view which has been abstracted from a concept along some dimension is what I call a
conceptual skeleton. In effect, we have dealt with conceptual skeletons all along, without
often using that name. For instance, many of the ideas concerning Bongard problems
could be rephrased using this terminology. It is always of interest, and possibly of
importance, when two or more ideas are discovered to share a conceptual skeleton. An
example is the bizarre set of concepts mentioned at the beginning of the Contrafactus: a
Bicyclops, a tandem unicycle, a teeter-teeter, the game of ping-ping, a one-way tie, a
two-sided Mobius strip, the "Bach twins", a piano concerto for two left hands, a one-
voice fugue, the act of clapping with one hand, a two-channel monaural phonograph, a
pair of eighth-backs. All of these ideas are "isomorphic" because they share this
conceptual skeleton:

                  a plural thing made singular and re-pluralized wrongly.

Two other ideas in this book which share that conceptual skeleton are (1) the Tortoise's
solution to Achilles' puzzle, asking for a word beginning and ending in "HE" (the
Tortoise's solution being the pronoun "HE", which collapses two occurrences into one),
and (2) the Pappus-Gelernter proof of the Pons As' norum Theorem, in which one triangle
is reperceived as two. Incidentally, these droll concoctions might be dubbed "demi-
doublets".
        A conceptual skeleton is like a set of constant features (as distinguished from
parameters or variables)-features which should not be slipped in a subjunctive instant
replay or mapping-operation. Having no parameters or variables of its own to vary, it can
be the invariant core of several different ideas. Each instance of it, such as "tandem
unicycle", does have layers of variability and so can be "slipped" in various ways.
        Although the name "conceptual skeleton" sounds absolute and rigid, actually
there is a lot of play in it. There can be conceptual skeletons on several different levels of
abstraction. For instance, the "isomorphism" between Bongard problems 70 and 71,
already pointed out, involves a higher-level conceptual skeleton than that needed to solve
either problem in isolation.




Artificial Intelligence:Prospects                                                         666

                              Multiple Representations
Not only must conceptual skeletons exist on different levels of abstraction; also, they
must exist along different conceptual dimensions. Let us take the following sentence as
an example:

         "The Vice President is the spare tire on the automobile of government."

How do we understand what it means (leaving aside its humor, which is of course a vital
aspect)? If you were told, "See our government as an automobile" without any prior
motivation, you might come up with any number of correspondences: steering wheel =
president, etc.. What are checks and balances? What are seat belts? Because the two
things being mapped are so different, it is almost inevitable that the mapping will involve
functional aspects. Therefore, you retrieve from your store of conceptual skeletons
representing parts of automobiles, only those having to do with function, rather than, say,
shape. Furthermore, it makes sense to work at a pretty high level of abstraction, where
"function" isn't taken in too narrow a context. Thus, of the two following definitions of
the function of a spare tire: (1) "replacement for a flat tire", and (2) "replacement for a
certain disabled part of a car", certainly the latter would be preferable, in this case. This
comes simply from the fact that an auto and a government are so different that they have
to be mapped at a high level of abstraction.
        Now when the particular sentence is examined, the mapping gets forced in one
respect-but it is not an awkward way, by any means. In fact, you already have a
conceptual skeleton for the Vice President, among many others, which says,
"replacement for a certain disabled part of government". Therefore the forced mapping
works comfortably. But suppose, for the sake of contrast, that you had retrieved another
conceptual skeleton for "spare tire"-say, one describing its physical aspects. Among other
things, it might say that a spare tire is "round and inflated". Clearly, this is not the right
way to go. (Or is it? As a friend of mine pointed out, some Vice Presidents are rather
portly, and most are quite inflated!)

                                     Ports of Access
One of the major characteristics of each idiosyncratic style of thought is how new
experiences get classified and stuffed into memory, for that defines the "handles" by
which they will later be retrievable. And for events, objects, ideas, and so on-for
everything that can be thought about-there is a wide variety of "handles". I am struck by
this each time I reach down to turn on my car radio, and find, to my dismay, that it is
already on! What has happened is that two independent representations are being used for
the radio. One is "music producer", the other is "boredom reliever". I am aware that the
music is on, but I am bored anyway, and before the two realizations have a chance to
interact, my reflex to reach




Artificial Intelligence:Prospects                                                         667

down has been triggered. The same reaching-down reflex one day occurred just after I'd
left the radio at a repair shop and was driving away, wanting to hear some music. Odd.
Many other representations for the same object exist, such as

               shiny silver-knob haver
               overheating-problems haver
               lying-on-my-back-over-hump-to-fix thing
               buzz-maker
               slipping-dials object
               multidimensional representation example

All of them can act as ports of access. Though they all are attached to my symbol for my
car radio, accessing that symbol through one does not open up all the others. Thus it is
unlikely that I will be inspired to remember lying on my back to fix the radio when I
reach down and turn it on. And conversely, when I'm lying on my back, unscrewing
screws, I probably won't think about the time I heard the Art of the Fugue on it. There are
"partitions" between these aspects of one symbol, partitions that prevent my thoughts
from spilling over sloppily, in the manner of free associations. My mental partitions are
important because they contain and channel the flow of my thoughts.
        One place where these partitions are quite rigid is in sealing off words for the
same thing in different languages. If the partitions were not strong, a bilingual person
would constantly slip back and forth between languages, which would be very
uncomfortable. Of course, adults learning two new languages at once often confuse
words in them. The partitions between these languages are flimsier, and can break down.
Interpreters are particularly interesting, since they can speak any of their languages as if
their partitions were inviolable and yet, on command, they can negate those partitions to
allow access to one language from the other, so they can translate. Steiner, who grew up
trilingual, devotes several pages in After Babel to the intermingling of French, English,
and German in the layers of his mind, and how his different languages afford different
ports of access onto concepts.

                                    Forced Matching
When two ideas are seen to share conceptual skeletons on some level of abstraction,
different things can happen. Usually the first stage is that you zoom in on both ideas, and,
using the higher-level match as a guide, you try to identify corresponding subideas.
Sometimes the match can be extended recursively downwards several levels, revealing a
profound isomorphism. Sometimes it stops earlier, revealing an analogy or similarity.
And then there are times when the high-level similarity is so compelling that, even if
there is no apparent lower-level continuation of the map, you just go ahead and make
one: this is the forced match.




Artificial Intelligence:Prospects                                                       668

         Forced matches occur every day in the political cartoons of newspapers: a
political figure is portrayed as an airplane, a boat, a fish, the Mona Lisa; a government is
a human, a bird, an oil rig; a treaty is a briefcase, a sword, a can of worms; on and on and
on. What is fascinating is how easily we can perform the suggested mapping, and to the
exact depth intended. We don't carry the mapping out too deeply or too shallowly.
         Another example of forcing one thing into the mold of another occurred when I
chose to describe the development of my Crab Canon in terms of meiosis. This happened
in stages. First, I noticed the common conceptual skeleton shared by the Crab Canon and
the image of chromosomes joined by a centromere; this provided the inspiration for the
forced match. Then I saw a high-level resemblance involving "growth", "stages", and
"recombination". Then I simply pushed the analogy as hard as I could. Tentativity-as in
the Bongard problem-solver-played a large role: I went forwards and backwards before
finding a match which I found appealing.
         A third example of conceptual mapping is provided by the Central Dogmap. I
initially noticed a high-level similarity between the discoveries of mathematical logicians
and those of molecular biologists, then pursued it on lower levels until I found a strong
analogy. To strengthen it further, I chose a Godel-numbering which imitated the Genetic
Code. This was the lone element of forced matching in the Central Dogmap.
         Forced matches, analogies, and metaphors cannot easily be separated out.
Sportscasters often use vivid imagery which is hard to pigeonhole. For instance, in a
metaphor such as "The Rams [football team are spinning their wheels", it is hard to say
just what image you are supposed to conjure up. Do you attach wheels to the team as a
whole\% Or to each player? Probably neither one. More likely, the image of wheels
spinning in mud or snow simply flashes before you for a brief instant, and then in some
mysterious way, just the relevant parts get lifted out and transferred to the team's
performance. How deeply are the football team and the car mapped onto each other in the
split second that you do this?

                                          Recap
Let me try to tie things together a little. I have presented a number of related ideas
connected with the creation, manipulation, and comparison of symbols. Most of them
have to do with slippage in some fashion, the idea being that concepts are composed of
some tight and some loose elements, coming from different levels of nested contexts
(frames). The loose ones can be dislodged and replaced rather easily, which, depending
on the circumstances, can create a "subjunctive instant replay", a forced match, or an
analogy. A fusion of two symbols may result from a process in which parts of each
symbol are dislodged and other parts remain.




Artificial Intelligence:Prospects                                                       669

                            Creativity and Randomness
It is obvious that we are talking about mechanization of creativity. But the this not a
contradiction in terms? Almost, but not really. Creativity s essence of that which is not
mechanical. Yet every creative act is mechanical-it has its explanation no less than a case
of the hiccups does. The mechanical substrate of creativity may be hidden from view, but
it exists. Conversely, there is something unmechanical in flexible programs, even today.
It may not constitute creativity, but when programs cease to be transparent to their
creators, then the approach to creativity has begun.
         It is a common notion that randomness is an indispensable ingredient of creative
acts. This may be true, but it does not have any bearing on the mechanizability-or rather,
programmability!-of creativity. The world is a giant heap of randomness; when you
mirror some of it inside your head, your head's interior absorbs a little of that
randomness. The triggering patterns of symbols, therefore, can lead you down the most
randomseeming paths, simply because they came from your interactions with a crazy,
random world. So it can be with a computer program, too. Randomness is an intrinsic
feature of thought, not something which has to be "artificially inseminated", whether
through dice, decaying nuclei, random number tables, or what-have-you. It is an insult to
human creativity to imply that it relies on such arbitrary sources.
         What we see as randomness is often simply an effect of looking at something
symmetric through a "skew" filter. An elegant example was provided by Salviati's two
ways of looking at the number it/4. Although the decimal expansion of 7r/4 is not literally
random, it is as random as one would need for most purposes: it is "pseudorandom".
Mathematics is full of pseudorandomness-plenty enough to supply all would-be creators
for all time.
         Just as science is permeated with "conceptual revolutions" on all levels at all
times, so the thinking of individuals is shot through and through with creative acts. They
are not just on the highest plane; they are everywhere. Most of them are small and have
been made a million times before-but they are close cousins to the most highly creative
and new acts. Computer programs today do not yet seem to produce many small
creations. Most of what they do is quite "mechanical" still. That just testifies to the fact
that they are not close to simulating the way we think-but they are getting closer.
         Perhaps what differentiates highly creative ideas from ordinary ones is some
combined sense of beauty, simplicity, and harmony. In fact, I have a favorite "meta-
analogy", in which I liken analogies to chords. The idea is simple: superficially similar
ideas are often not deeply related; and deeply related ideas are often superficially
disparate. The analogy to chords is natural: physically close notes are harmonically
distant (e.g., E-F-G); and harmonically close notes are physically distant (e.g., G-E-B).
Ideas that share a conceptual skeleton resonate in a sort of conceptual analogue to
harmony; these harmonious "idea-chords" are often widely separated, as




Artificial Intelligence:Prospects                                                       670

measured on an imaginary "keyboard of concepts". Of course, it doesn't suffice to reach
wide and plunk down any old way-you may hit a seventh or a ninth! Perhaps the present
analogy is like a ninth-chord-wide but dissonant.

                         Picking up Patterns on All Levels
Bongard problems were chosen as a focus in this Chapter because when you study them,
you realize that the elusive sense for patterns which we humans inherit from our genes
involves all the mechanisms of representation of knowledge, including nested contexts,
conceptual skeletons and conceptual mapping, slippability, descriptions and meta-
descriptions and their interactions, fission and fusion of symbols, multiple representations
(along different dimensions and different levels of abstraction), default expectations, and
more.
        These days, it is a safe bet that if some program can pick up patterns in one area,
it will miss patterns in another area which, to us, are equally obvious. You may remember
that I mentioned this back in Chapter 1, saying that machines can be oblivious to
repetition, whereas people cannot. For instance, consider SHRDLU. If Eta Oin typed the
sentence "Pick up a big red block and put it down" over and over again, SHRDLU would
cheerfully react in the same way over and over again, exactly as an adding machine will
print out "4" over and over again, if a human being has the patience to type "2+2" over
and over again. Humans aren't like that; if some pattern occurs over and over again, they
will pick it up. SHRDLU wasn't built with the potential for forming new concepts or
recognizing patterns: it had no sense of over and overview.

                            The Flexibility of Language
SHRDLU's language-handling capability is immensely flexible-within limits. SHRDLU
can figure out sentences of great syntactical complexity, or sentences with semantic
ambiguities as long as-they can- be resolved by inspecting the data base-but it cannot
handle "hazy" language. For instance, consider the sentence "How many blocks go on top
of each other to make a steeple?" We understand it immediately, yet it does not make
sense if interpreted literally. Nor is it that some idiomatic phrase has been used. "To go
on top of each other" is an imprecise phrase which nonetheless gets the desired image
across quite well to a human. Few people would be misled into visualizing a paradoxical
setup with two blocks each of which is on top of the other-or blocks which are "going"
somewhere or other.
        The amazing thing about language is how imprecisely we use it and still manage
to get away with it. SHRDLU uses words in a "metallic" way, while people use them in a
"spongy" or "rubbery" or even "Nutty-Puttyish" way. If words were nuts and bolts,
people could make any bolt fit into any nut: they'd just squish the one into the other, as in
some surrealistic




Artificial Intelligence:Prospects                                                        671

painting where everything goes soft. Language, in human hands, becomes almost like a
fluid, despite, the coarse grain of its components.
        Recently, Al research in natural language understanding has turned away
somewhat from the understanding of single sentences in isolation, and more towards
areas such as understanding simple children's stories. Here is a well-known children's
joke which illustrates the open-endedness of real-life situations:

    A man took a ride in an airplane.
       Unfortunately, he fell out.
    Fortunately, he had a parachute on.
       Unfortunately, it didn't work.
    Fortunately, there was a haystack below him.
       Unfortunately, there was a pitchfork sticking out of it.
    Fortunately, he missed the pitchfork.
       Unfortunately, he missed the haystack.

It can be extended indefinitely. To represent this silly story in a frame-based system
would be extremely complex, involving jointly activating frames for the concepts of man,
airplane, exit, parachute, falling, etc., etc.

                              Intelligence and Emotions
Or consider this tiny yet poignant story:

    Margie was holding tightly to the string of her beautiful new balloon. Suddenly, a
    gust of wind caught it. The wind carried it into a tree. The balloon hit a branch
    and burst. Margie cried and cried.'

To understand this story, one needs to read many things between the lines. For instance:
Margie is a little girl. This is a toy balloon with a string for a child to hold. It may not be
beautiful to an adult, but in a child's eye, it is. She is outside. The "it" that the wind
caught was the balloon. The wind did not pull Margie along with the balloon; Margie let
go. Balloons can break on contact with any sharp point. Once they are broken, they are
gone forever. Little children love balloons and can be bitterly disappointed when they
break. Margie saw that her balloon was broken. Children cry when they are sad. "To cry
and cry" is to cry very long and hard. Margie cried and cried because of her sadness at
her balloon's breaking.
        This is probably only a small fraction of what is lacking at the surface level. A
program must have all this knowledge in order to get at what is going on. And you might
object that, even if it "understands" in some intellectual sense what has been said, it will
never really understand, until it, too, has cried and cried. And when will a computer do
that? This is the kind of humanistic point which Joseph Weizenbaum is concerned with
making in his book Computer Power and Human Reason, and I think it is an important
issue; in fact, a very, very deep issue. Unfortunately, many Al workers at this time are
unwilling, for various reasons, to take this sort of point




Artificial Intelligence:Prospects                                                          672

seriously. taut in some ways, those Al workers are right: it is a little premature to think
about computers crying; we must first think about rules for computers to deal with
language and other things; in time, we'll find ourselves face to face with the deeper
issues.

                                    AI Has Far to Go
Sometimes it seems that there is such a complete absence of rule-governed behavior that
human beings just aren't rule-governed. But this is an illusion-a little like thinking that
crystals and metals emerge from rigid underlying laws, but that fluids or flowers don't.
We'll come back to this question in the next Chapter.

   The process of logic itself working internally in the brain may be more analogous
   to a succession of operations with symbolic pictures, a sort of abstract analogue of
   the Chinese alphabet or some Mayan description of events-except that the
   elements are not merely words but more like sentences or whole stories with
   linkages between them forming a sort of meta- or super-logic with its own rules.'

        It is hard for most specialists to express vividly-perhaps even to remember-what
originally sparked them to enter their field. Conversely, someone on the outside may
understand a field's special romance and may be able to articulate it precisely. I think that
is why this quote from Ulam has appeal for me, because it poetically conveys the
strangeness of the enterprise of Al, and yet shows faith in it. And one must run on faith at
this point, for there is so far to go!

                         Ten Questions and Speculations
To conclude this Chapter, I would like to present ten "Questions and Speculations" about
Al. I would not make so bold as to call them "Answers"-these are my personal opinions.
They may well change in some ways, as I learn more and as Al develops more. (In what
follows, the term "Al program" means a program which is far ahead of today's programs;
it means an "Actually Intelligent" program. Also, the words "program" and "computer"
probably carry overly mechanistic connotations, but let us stick with them anyway.)

Question: Will a computer program ever write beautiful music?
 Speculation: Yes, but not soon. Music is a language of emotions, and until programs
    have emotions as complex as ours, there is no way a program will write anything
    beautiful. There can be "forgeries” shallow imitations of the syntax of earlier music-
    but despite what one might think at first, there is much more to musical expression
    than can be captured in syntactical rules. There will be no new kinds of beauty
    turned up for a long time by computer music-composing programs. Let me carry this
    thought a little further. To think-and I have heard this suggested-that we might soon
    be able to command a preprogrammed mass-produced mail-order twenty-dollar
    desk-model "music box" to bring forth from its sterile circuitry pieces which Chopin
    or Bach might have written had they lived longer is a grotesque and shameful
    misestimation of the depth of the human spirit. A "program" which could produce


Artificial Intelligence:Prospects                                                        673

    music as they did would have to wander around the world on its own, fighting its
    way through the maze of life and feeling every moment of it. It would have to
    understand the joy and loneliness of a chilly night wind, the longing for a cherished
    hand, the inaccessibility of a distant town, the heartbreak and regeneration after a
    human death. It would have to have known resignation and worldweariness, grief
    and despair, determination and victory, piety and awe. In it would have had to
    commingle such opposites as hope and fear, anguish and jubilation, serenity and
    suspense. Part and parcel of it would have to be a sense of grace, humor, rhythm, a
    sense of the unexpected-and of course an exquisite awareness of the magic of fresh
    creation. Therein, and therein only, lie the sources of meaning in music.

Question: Will emotions be explicitly programmed into a machine?
  Speculation: No. That is ridiculous. Any direct simulation of emotions-PARRY, for
     example-cannot approach the complexity of human emotions, which arise
     indirectly from the organization of our minds. Programs or machines will acquire
     emotions in the same way: as by-products of their structure, of the way in which
     they are organized-not by direct programming. Thus, for example, nobody will
     write a "falling-in-love" subroutine, any more than they would write a "mistake-
     making" subroutine. "Falling in love" is a description which we attach to a complex
     process of a complex system; there need be no single module inside the system
     which is solely responsible for it, however!

Question: Will a thinking computer be able to add fast?
  Speculation: Perhaps not. We ourselves are composed of hardware which does fancy
     calculations but that doesn't mean that our symbol level, where "we" are, knows
     how to carry out the same fancy calculations. Let me put it this way: there's no way
     that you can load numbers into your own neurons to add up your grocery bill.
     Luckily for you, your symbol level (i.e., you) can't gain access to the neurons which
     are doing your thinking-otherwise you'd get addle-brained. To paraphrase Descartes
     again:

                            "I think; therefore I have no access
                                 to the level where I sum."

   Why should it not be the same for an intelligent program? It mustn't be allowed to gain
    access to the circuits which are doing its thinking otherwise it'll get addle-CPU'd.
    Quite seriously, a machine that can pass the Turing test may well add as slowly as
    you or I do, and for




Artificial Intelligence:Prospects                                                     674

   similar reasons. It will represent the number 2 not just by the two bits "10", but as a
      full-fledged concept the way we do, replete with associations such as its homonyms
      "too" and "to", the words "couple" and "deuce", a host of mental images such as
      dots on dominos, the shape of the numeral '2', the notions of alternation, evenness,
      oddness, and on and on ... With all this "extra baggage" to carry around, an
      intelligent program will become quite slothful in its adding. Of course, we could
      give it a ' pocket calculator , so to speak (or build one in). Then it could answer
      very fast, but its performance would be just like that of a person with a pocket
      calculator. There would be two separate parts to the machine: a reliable but
      mindless part and an intelligent but fallible part. You couldn't rely on the composite
      system to be reliable, any more than a composite of person and machine is
      necessarily reliable. So if it's right answers you're after, better stick to the pocket
      calculator alone-don't throw in the intelligence!

Question: Will there be chess programs that can beat anyone?
  Speculation: No. There may be programs which can beat anyone at chess, but they
     will not be exclusively chess players. They will be programs of general
     intelligence, and they will be just as temperamental as people. "Do you want to play
     chess?" "No, I'm bored with chess. Let's talk about poetry." That may be the kind of
     dialogue you could have with a program that could beat everyone. That is because
     real intelligence inevitably depends on a total overview capacity-that is, a
     programmed ability to "jump out of the system", so to speak-at least roughly to the
     extent that we have that ability. Once that is present, you can't contain the program;
     it's gone beyond that certain critical point, and you just have to face the facts of
     what you've wrought.

Question: Will there be special locations in memory which store parameters governing
  the behavior of the program, such that if you reached in and changed them, you would
  be able to make the program smarter or stupider or more creative or more interested in
  baseball? In short, would you be able to "tune" the program by fiddling with it on a
  relatively low level?
  Speculation: No. It would be quite oblivious to changes of any particular elements in
     memory, just as we stay almost exactly the same though thousands of our neurons
     die every day(!). If you fuss around too heavily, though, you'll damage it, just as if
     you irresponsibly did neurosurgery on a human being. There will be no "magic"
     location in memory where, for instance, the "IQ" of the program sits. Again, that
     will be a feature which emerges as a consequence of lower-level behavior, and
     nowhere will it sit explicitly. The same goes for such things as "the number of
     items it can hold in short-term memory", "the amount it likes physics", etc., etc.

Question: Could you "tune" an Al program to act like me, or like you-or halfway between
us?




Artificial Intelligence:Prospects                                                        675

   Speculation: No. An intelligent program will not be chameleon-like, any more than
     people are. ,It will rely on the constancy of its memories, and will not be able to flit
     between personalities. The idea of changing internal parameters to "tune to a new
     personality" reveals a ridiculous underestimation of the complexity of personality.

Question: Will there be a "heart" to an Al program, or will it simply consist of "senseless
  loops and sequences of trivial operations" (in the words of Marvin Minskys)?
  Speculation: If we could see all the way to the bottom, as we can a shallow pond, we
     would surely see only "senseless loops and sequences of trivial operations"-and we
     would surely not see any "heart". Now there are two kinds of extremist views on
     AI: one says that the human mind is, for fundamental and mysterious reasons,
     unprogrammable. The other says that you merely need to assemble the appropriate
     "heuristic devices-multiple optimizers, pattern-recognition tricks, planning
     algebras, recursive administration procedures, and the like",' and you will have
     intelligence. I find myself somewhere in between, believing that the "pond" of an
     Al program will turn out to be so deep and murky that we won't be able to peer all
     the way to the bottom. If we look from the top, the loops will be invisible, just as
     nowadays the current-carrying electrons are invisible to most programmers. When
     we create a program that passes the Turing test, we will see a "heart" even though
     we know it's not there.

Question: Will Al programs ever become "superintelligent"?
  Speculation: I don't know. It is not clear that we would be able to understand or relate
     to a "superintelligence", or that the concept even makes sense. For instance, our
     own intelligence is tied in with our speed of thought. If our reflexes had been ten
     times faster or slower, we might have developed an entirely different set of
     concepts with which to describe the world. A creature with a radically different
     view of the world may simply not have many points of contact with us. I have often
     wondered if there could be, for instance, pieces of music which are to Bach as Bach
     is to folk tunes: "Bach squared", so to speak. And would I be able to understand
     them? Maybe there is such music around me already, and I just don't recognize it,
     just as dogs don't understand language. The idea of superintelligence is very
     strange. In any case, I don't think of it as the aim of Al research, although if we ever
     do reach the level of human intelligence, superintelligence will undoubtedly be the
     next goal-not only for us, but for our Al-program colleagues, too, who will be
     equally curious about Al and superintelligence. It seems quite likely that Al
     programs will be extremely curious about Al in general-understandably.

Question: You seem to be saying that AI programs will be virtually identical to people,
  then. Won't there be any differences?




Artificial Intelligence:Prospects                                                        676

   Speculation: Probably the differences between Al programs and people will be larger
     than the differences between most people. It is almost impossible to imagine that
     the "body" in which an Al program is housed would not affect it deeply. So unless
     it had an amazingly faithful replica of a human body-and why should it?-it would
     probably have enormously different perspectives on what is important, what is
     interesting, etc. Wittgenstein once made the amusing comment, "If a lion could
     speak, we would not understand him." It makes me think of Rousseau's painting of
     the gentle lion and the sleeping gypsy on the moonlit desert. But how does
     Wittgenstein know? My guess is that any Al program would, if comprehensible to
     us, seem pretty alien. For that reason, we will have a very hard time deciding when
     and if we really are dealing with an Al program, or just a "weird" program.

Question: Will we understand what intelligence and consciousness and free will and "I"
  are when we have made an intelligent program?
  Speculation: Sort of-it all depends on what you mean by "understand". On a gut level,
     each of us probably has about as good an understanding as is possible of those
     things, to start with. It is like listening to music. Do you really understand Bach
     because you have taken him apart? Or did you understand it that time you felt the
     exhilaration in every nerve in your body? Do we understand how the speed of light
     is constant in every inertial reference frame? We can do the math, but no one in the
     world has a truly relativistic intuition. And probably no one will ever understand
     the mysteries of intelligence and consciousness in an intuitive way. Each of us can
     understand people, and that is probably about as close as you can come.




Artificial Intelligence:Prospects                                                    677

